\documentclass[a4paper,11pt]{article}
\usepackage[utf8]{inputenc}
\usepackage[french]{babel}
\usepackage[T1]{fontenc}
\usepackage{amsmath,amssymb,amsthm}
\usepackage{geometry}
\geometry{margin=2cm}
\usepackage{enumitem}
\usepackage{hyperref}
\usepackage{booktabs}
\usepackage{array}
\usepackage{xcolor}
\usepackage{listings}
\usepackage{titlesec}
\usepackage{fancyhdr}
\usepackage{lastpage}
\usepackage{graphicx}
\usepackage{etoolbox}
\AtBeginDocument{\shorthandoff{;:!?}}

% ---- Style des titres ----
\titleformat{\section}{\large\bfseries\color{blue!60!black}}{}{0pt}{}[\vspace{-0.5ex}\rule{\textwidth}{0.4pt}]
\titleformat{\subsection}{\bfseries\color{blue!40!black}}{}{0pt}{}
\titleformat{\subsubsection}{\itshape\color{blue!30!black}}{}{0pt}{}

% ---- Style listings ----
\lstdefinestyle{monstyle}{
  frame=single,
  numbers=left,
  numbersep=5pt,
  breaklines=true,
  basicstyle=\small\ttfamily,
  backgroundcolor=\color{gray!5},
  rulecolor=\color{gray!40},
  columns=fullflexible,
  keepspaces=true,
  showstringspaces=false
}
\lstset{style=monstyle}

% ---- En-têtes / pieds ----
\pagestyle{fancy}
\fancyhf{}
\fancyhead[L]{\small STA211 -- Données manquantes}
\fancyhead[R]{\small Fiche récapitulative}
\fancyfoot[C]{\thepage/\pageref{LastPage}}
\renewcommand{\headrulewidth}{0.4pt}

% ---- Theorem-like environment ----
\theoremstyle{definition}
\newtheorem*{definition}{Définition}
\theoremstyle{remark}
\newtheorem*{remark}{Remarque}
\newtheorem*{important}{À retenir}

\title{\textbf{Fiche récapitulative -- STA211}\\[4pt]
\large Gestion des données manquantes}
\author{}
\date{10 juin 2026}

\begin{document}
\maketitle
\thispagestyle{empty}

\begin{abstract}
Cette fiche synthétise le cours sur la gestion des données manquantes
(V. Audigier, N. Niang). Les données manquantes sont omniprésentes en
data-mining~: questionnaires incomplets, fusion de fichiers, défaillances
de saisie, etc. Leur traitement est indispensable car la plupart des
méthodes d'analyse ne sont pas conçues pour les données incomplètes.
\end{abstract}

\vfill
\tableofcontents
\newpage

% ===================================================================
\section{Introduction}
% ===================================================================

Les données manquantes constituent un problème incontournable dans
l'application des techniques de data-mining. L'augmentation de la taille
des données fait mécaniquement augmenter la probabilité d'en rencontrer.

Les causes sont multiples~:
\begin{itemize}
  \item Oubli de saisie, questionnaires partiellement remplis
  \item Refus de répondre (questions sensibles~: revenus, consommation,
    etc.)
  \item Fusion de jeux de données~: concaténer des fichiers fait
    apparaître des blocs verticaux de données manquantes car les
    variables diffèrent.
\end{itemize}

\textbf{Deux approches}~:
\begin{enumerate}
  \item Pré-traitement pour se ramener à des données complètes (imputation)
  \item Méthodes avancées qui s'appliquent directement sur les données
    incomplètes (EM, pondération)
\end{enumerate}

La méthode naïve consistant à supprimer les individus incomplets
(\textit{liste-wise deletion}) n'est pas satisfaisante~:
\begin{itemize}
  \item Les individus complets ne sont pas nécessairement représentatifs
    $\rightarrow$ biais.
  \item Le nombre d'individus complets devient vite très petit quand le
    nombre de variables est grand.
\end{itemize}
Exemple~: si 5\% des valeurs sont manquantes pour chaque variable (MCAR),
un jeu de 50 variables contient en moyenne seulement 8\% d'individus
complets.

% ===================================================================
\section{Mécanismes de manquance}
% ===================================================================

On note $\mathbf{X}_{n\times p}$ le tableau de données complet (fictif)
et $\mathbf{R}_{n\times p}$ la matrice d'indication de manquance~:
$r_{ij} = 1$ si $x_{ij}$ est observé, $0$ sinon.

Les méthodes employées pour gérer les données manquantes sont conditionnées
par le type de mécanisme.

\subsection{MCAR (Missing Completely At Random)}

La probabilité de manquance est indépendante de toutes les données
(observées et manquantes)~:
\[
P(R \mid X) = P(R).
\]

\textbf{Propriété}~: les individus complets forment un sous-échantillon
représentatif. L'analyse sur les cas complets n'est pas biaisée, mais
la variance des estimateurs est augmentée.

\textbf{Exemple}~: questionnaire perdu aléatoirement.

\subsection{MAR (Missing At Random)}

La probabilité de manquance dépend des données observées mais pas des
données manquantes~:
\[
P(R \mid X) = P(R \mid X^{\text{obs}}).
\]

\textbf{Propriété}~: l'analyse par cas complets est biaisée. MAR est
l'hypothèse par défaut de la plupart des méthodes d'imputation.

\textbf{Exemple}~: enquête de satisfaction~: seules les personnes
insatisfaites en janvier sont réinterrogées en février. La manquance
en février dépend de la réponse observée en janvier.

\begin{important}
  L'hypothèse MAR est non testable. Il est recommandé d'inclure des
  \textbf{variables auxiliaires} (peu/pas de manquantes, sans intérêt
  scientifique propre) qui aident à expliquer la manquance et rendent
  MAR plus crédible.
\end{important}

\subsection{MNAR (Missing Not At Random)}

La probabilité de manquance dépend des données manquantes elles-mêmes~:
\[
P(R \mid X) = P(R \mid X^{\text{obs}}, X^{\text{mis}}).
\]

\textbf{Exemple}~: les personnes à haut revenu ne déclarent pas leur
revenu.

\textbf{Traitement}~: méthodes spécifiques (pattern mixture models,
selection models), moins développées et plus complexes que pour MAR.

% ===================================================================
\section{Identification des mécanismes}
% ===================================================================

Bien qu'il soit impossible de \emph{démontrer} le mécanisme, l'analyse
exploratoire aide à l'identifier~:

\begin{enumerate}
  \item \textbf{Analyse univariée du dispositif} $\mathbf{R}$~:
    proportion de manquantes par variable.
  \item \textbf{Analyse bivariée}~: lien entre la manquance d'une
    variable et les valeurs observées d'autres variables.
    Exemple~: distribution de \textit{Credit Amount} selon que
    \textit{Duration} est observée ou non.
  \item \textbf{Analyse multivariée}~: régression logistique de la
    manquance sur les variables observées.
\end{enumerate}

\begin{remark}
  Ces analyses ne permettent pas de trancher entre MAR et MNAR (tous
  deux peuvent produire des différences de distribution). Seule la
  connaissance du mécanisme de production des données permet de
  décider.
\end{remark}

% ===================================================================
\section{Traitement par imputation simple}
% ===================================================================

L'imputation simple remplace chaque valeur manquante par une valeur
plausible (une seule valeur par cellule).

\subsection{Typologie des méthodes}

\paragraph{Imputation marginale (univariée)}~:
moyenne, médiane, tirage aléatoire dans les valeurs observées.
\textbf{Inconvénient}~: détruit les relations entre variables.

\paragraph{Imputation par régression}~:
\begin{enumerate}
  \item Ajuster un modèle de régression sur les données observées.
  \item Prédire les valeurs manquantes à partir de ce modèle.
  \item Ajouter une perturbation aléatoire (\textbf{imputation
    stochastique}) pour mieux respecter la distribution.
\end{enumerate}

\paragraph{Imputation par $k$ plus proches voisins ($k$-NN)}~:
\begin{enumerate}
  \item Identifier les $k$ individus complets les plus proches de
    l'individu à imputer (distance euclidienne).
  \item Tirer une valeur parmi leurs $k$ valeurs observées (ou
    faire la moyenne).
\end{enumerate}

\paragraph{Choix du modèle}~:
\begin{quote}
  «~Choisir des modèles d'imputation au moins aussi complexes que la
  méthode qui sera appliquée sur les données imputées~» (Schafer, 2003).
\end{quote}
Si l'analyse finale comporte des interactions, l'imputation doit les
préserver.

\subsection{Imputation séquentielle (plusieurs variables)}

Algorithme~:
\begin{enumerate}
  \item Ordonner les variables de la moins manquante à la plus manquante.
  \item Initialiser les valeurs manquantes (tirage dans les observées).
  \item Pour chaque variable $j$ (dans l'ordre)~:
    \begin{itemize}
      \item Ajuster le modèle d'imputation de $X_j$ sur les autres
        variables (données observées de $X_j$ uniquement).
      \item Mettre à jour les valeurs imputées de $X_j$ avec ce modèle.
    \end{itemize}
  \item Répéter l'étape 3 jusqu'à convergence (stabilisation de la
    distribution imputée).
\end{enumerate}
Les logiciels proposent des modèles par défaut selon la nature des
variables (régression linéaire pour quantitative, logistique pour
binaire, etc.).

\subsection{Validation du modèle d'imputation}

Comparer les distributions des données imputées et observées~:
\begin{itemize}
  \item Une différence de \textbf{tendance centrale} est attendue sous MAR.
  \item Une différence de \textbf{forme} indique un modèle mal adapté.
\end{itemize}
Outils~: histogrammes superposés, boxplots, analyses bivariées/multivariées.

\subsection{Modèles joints}

Alternative aux approches séquentielles~: spécifier un modèle pour la
distribution \textbf{jointe} de toutes les variables.

\begin{itemize}
  \item \textbf{Données quantitatives}~: modèle gaussien multivarié
    (package \texttt{Amelia}).
  \item \textbf{Données qualitatives}~: modèle log-linéaire, classes
    latentes (package \texttt{cat}).
  \item \textbf{Données mixtes}~: \textit{general location model}
    (package \texttt{mix}).
  \item \textbf{Imputation par analyse factorielle}~: ACP, ACM, AFDM
    (package \texttt{missMDA}). Particulièrement adapté à la grande
    dimension, la multicolinéarité et les données mixtes.
\end{itemize}

Avantage~: rapidité (toutes les variables imputées simultanément).
Inconvénient~: moins flexible que l'approche séquentielle.

% ===================================================================
\section{Imputation multiple -- MICE}
% ===================================================================

L'imputation simple ignore l'incertitude sur les valeurs imputées~:
les intervalles de confiance construits après imputation simple ont un
taux de couverture trop faible.

L'\textbf{imputation multiple} génère $M$ tableaux imputés (typiquement
$M \ge 3$, souvent $M=5$ ou $M=20$)~:

\begin{enumerate}
  \item Générer $M$ jeux de paramètres du modèle d'imputation (par
    bootstrap ou tirage dans la distribution a posteriori).
  \item Pour chaque jeu, imputer les données $\rightarrow$ $M$ tableaux
    complets.
  \item Appliquer la méthode d'analyse sur chaque tableau.
  \item Agréger les $M$ résultats avec les \textbf{règles de Rubin}.
\end{enumerate}

\subsection{Règles de Rubin}

Soit $\hat{\theta}_m$ l'estimation du paramètre d'intérêt sur le
$m$-ième tableau imputé, et $\widehat{\operatorname{Var}}(\hat{\theta}_m)$
sa variance estimée.

\textbf{Estimation agrégée}~:
\[
\hat{\theta} = \frac{1}{M}\sum_{m=1}^{M} \hat{\theta}_m.
\]

\textbf{Variance totale}~:
\[
\widehat{\operatorname{Var}}(\hat{\theta})
= \frac{1}{M}\sum_{m=1}^{M} \widehat{\operatorname{Var}}(\hat{\theta}_m)
+ \left(1 + \frac{1}{M}\right)
  \frac{1}{M-1}\sum_{m=1}^{M} (\hat{\theta}_m - \hat{\theta})^2.
\]

Le premier terme est la \textbf{variance intra-imputation} (incertitude
dans chaque échantillon), le second la \textbf{variance inter-imputation}
(incertitude due aux données manquantes).

\begin{remark}
  Les règles de Rubin s'appliquent essentiellement aux modèles linéaires
  généralisés (régression linéaire, logistique, ANOVA). Pour d'autres
  méthodes, l'agrégation peut être difficile.
\end{remark}

\subsection{Overimputation}

Technique de validation~: générer des valeurs plausibles également pour
les \textbf{données observées}. En générant plusieurs centaines de
tableaux, on construit un intervalle de prédiction à 95~\% pour chaque
valeur observée. Si le modèle est bien choisi, 95~\% des valeurs
observées sont contenues dans leur intervalle.

% ===================================================================
\section{Autres méthodes de traitement}
% ===================================================================

\subsection{Algorithmes EM et DA}

\begin{itemize}
  \item \textbf{EM (Expectation-Maximization)}~: alterne imputation
    (espérance) et estimation des paramètres par maximum de vraisemblance.
  \item \textbf{DA (Data Augmentation)}~: version bayésienne qui alterne
    génération des données imputées et tirage des paramètres dans leur
    distribution a posteriori.
\end{itemize}

\textbf{Différence avec l'imputation}~: le modèle d'imputation est le
même que le modèle d'analyse. On ne peut pas, par exemple, imputer par
$k$-NN et analyser par régression linéaire (contrairement à l'imputation).

\subsection{Méthodes par pondération}

Attribuer un poids aux individus complets, inversement proportionnel à
leur probabilité d'être manquants. Issu de la théorie des sondages.

\textbf{Avantage}~: pas besoin de modéliser la distribution des données.
\textbf{Inconvénient}~: nécessite de modéliser le mécanisme de manquance
(estimation difficile des poids en pratique).

% ===================================================================
\section{Code R -- Exemples}
% ===================================================================

\begin{lstlisting}[language=R, caption=Imputation multiple avec mice]
library(mice)

# Chargement et inspection
data(nhanes)
md.pattern(nhanes)

# Imputation multiple
imp <- mice(nhanes, m = 5, method = 'pmm',
            seed = 123, printFlag = FALSE)

# Convergence
plot(imp)

# Comparaison distributions
densityplot(imp)

# Analyse et combinaison
fit <- with(imp, lm(bmi ~ age + chl))
pooled <- pool(fit)
summary(pooled)
\end{lstlisting}

\begin{lstlisting}[language=R, caption=Imputation par forêts aléatoires (randomForest)]
library(mice)
imp_rf <- mice(nhanes, m = 5,
               method = 'rf',  # random forest
               seed = 123, printFlag = FALSE)
fit_rf <- with(imp_rf, lm(bmi ~ age + chl))
summary(pool(fit_rf))
\end{lstlisting}

\begin{lstlisting}[language=R, caption=Imputation par analyse factorielle (missMDA)]
library(missMDA)
data(ozone)

# Estimation du nombre de composantes
nb <- estim_ncpPCA(ozone[, -1], ncp.max = 5)

# Imputation par ACP
res.impute <- imputePCA(ozone[, -1], ncp = nb$ncp)
head(res.impute$completeObs)
\end{lstlisting}

\begin{lstlisting}[language=R, caption=Imputation avec Amelia (modèle gaussien)]
library(Amelia)
data(freetrade)
amelia_out <- amelia(freetrade, m = 5,
                     ts = 'year', cs = 'country')
summary(amelia_out)
\end{lstlisting}

\begin{lstlisting}[language=R, caption=Codes pour simuler l'impact de l'imputation simple]
set.seed(123)
nsim <- 2000
res <- matrix(NA, nsim, 2)

for (ii in 1:nsim) {
  x <- rnorm(500)
  y <- 2 * x + rnorm(500, sd = 0.5)

  # Ajout de donnees manquantes MAR
  prob <- pnorm(1.2 * x - 0.5)
  ismar <- sapply(prob,
    FUN = function(p) sample(c(TRUE, FALSE), 1,
                             prob = c(p, 1 - p)))
  ymar <- y; ymar[ismar] <- NA

  # Imputation simple par regression stochastique
  res.lm <- lm(ymar[!ismar] ~ x[!ismar])
  ystoch <- ymar
  ystoch[ismar] <- cbind(1, x[ismar]) %*%
    res.lm$coefficients +
    rnorm(sum(ismar), sd = summary(res.lm)$sigma)

  # Intervalle de confiance
  se <- sd(ystoch) / sqrt(500)
  ci_inf <- mean(ystoch) - qt(0.975, 499) * se
  ci_sup <- mean(ystoch) + qt(0.975, 499) * se
  res[ii, ] <- c(mean(ystoch), 0 >= ci_inf & 0 <= ci_sup)
}

# Taux de couverture (devrait etre < 0.95)
cat("Taux de couverture moyen:",
    mean(res[, 2]), "\n")
\end{lstlisting}

% ===================================================================
\section*{Questions clés (QCM)}
% ===================================================================

\addcontentsline{toc}{section}{Questions clés (QCM)}

\begin{enumerate}

  \item Le mécanisme MAR signifie que la probabilité de manquance~:
    \begin{enumerate}
      \item Est indépendante de toutes les données
      \item Dépend des données observées uniquement
      \item Dépend des données manquantes uniquement
      \item Dépend des données observées et manquantes
    \end{enumerate}
    \textbf{Réponse~: b}

  \item Sous un mécanisme MCAR~:
    \begin{enumerate}
      \item L'analyse sur les cas complets est biaisée
      \item Les individus complets sont un sous-échantillon
        représentatif
      \item Il faut obligatoirement imputer les données
      \item La variance des estimateurs est plus faible
    \end{enumerate}
    \textbf{Réponse~: b}

  \item L'hypothèse MAR est~:
    \begin{enumerate}
      \item Testable par un test statistique
      \item Non testable
      \item Toujours vérifiée en pratique
      \item Plus restrictive que MCAR
    \end{enumerate}
    \textbf{Réponse~: b}

  \item Dans MICE, que signifie l'acronyme~?
    \begin{enumerate}
      \item Multiple Imputation by Chained Equations
      \item Missing Information Criterion Estimate
      \item Multivariate Imputation for Complete Estimation
      \item Model-based Imputation for Censored Events
    \end{enumerate}
    \textbf{Réponse~: a}

  \item Quelle est la conséquence de l'imputation simple sur les
    intervalles de confiance~?
    \begin{enumerate}
      \item Ils sont trop larges
      \item Ils sont trop courts (sous-couverture)
      \item Ils ne sont pas affectés
      \item Ils deviennent exacts
    \end{enumerate}
    \textbf{Réponse~: b}

  \item Dans les règles de Rubin, la variance totale $T$ est~:
    \begin{enumerate}
      \item $T = \bar{U}$ (variance intra uniquement)
      \item $T = B$ (variance inter uniquement)
      \item $T = \bar{U} + (1+1/M)B$
      \item $T = \bar{U} - (1+1/M)B$
    \end{enumerate}
    \textbf{Réponse~: c}

  \item Quel package R propose l'imputation multiple par modèle
    gaussien joint~?
    \begin{enumerate}
      \item \texttt{mice}
      \item \texttt{Amelia}
      \item \texttt{VIM}
      \item \texttt{missMDA}
    \end{enumerate}
    \textbf{Réponse~: b}

  \item L'approche par analyse factorielle pour l'imputation
    (package \texttt{missMDA}) est particulièrement adaptée~:
    \begin{enumerate}
      \item Aux petites données
      \\item À la grande dimension et la multicolinéarité
      \item Aux données qualitatives uniquement
      \item Aux séries temporelles
    \end{enumerate}
    \textbf{Réponse~: b}

  \item La technique d'\textit{overimputation} consiste à~:
    \begin{enumerate}
      \item Imputer plus de variables que nécessaire
      \item Générer des valeurs plausibles pour les données observées
      \item Supprimer les valeurs imputées
      \item Augmenter le nombre $M$ d'imputations
    \end{enumerate}
    \textbf{Réponse~: b}

  \item Quelle est la différence entre imputation et algorithme EM~?
    \begin{enumerate}
      \item Aucune, ce sont les mêmes méthodes
      \item En imputation, le modèle d'imputation peut différer du
        modèle d'analyse
      \item L'EM ne nécessite pas de modèle
      \item L'imputation ne fonctionne que sous MCAR
    \end{enumerate}
    \textbf{Réponse~: b}

\end{enumerate}

% ===================================================================
\begin{thebibliography}{9}
% ===================================================================

\bibitem{little2002}
R.J.A.~Little, D.B.~Rubin.
\newblock \emph{Statistical Analysis with Missing Data}.
\newblock 2nd ed., Wiley, 2002.

\bibitem{vanbuuren2012}
S.~van~Buuren.
\newblock \emph{Flexible Imputation of Missing Data}.
\newblock CRC Press, 2012.

\bibitem{audigier2026}
V.~Audigier, N.~Niang.
\newblock STA211~: gestion des données manquantes.
\newblock Cours CNAM, 2026.

\bibitem{schafer2003}
J.L.~Schafer.
\newblock Multiple imputation in multivariate problems when the imputation
  and analysis models differ.
\newblock \emph{Statistica Neerlandica}, 57(1):19--35, 2003.

\bibitem{josse2016}
J.~Josse, F.~Husson.
\newblock missMDA~: a package for handling missing values in multivariate
  data analysis.
\newblock \emph{Journal of Statistical Software}, 70(1):1--31, 2016.

\end{thebibliography}

\end{document}
